% Options for packages loaded elsewhere
\PassOptionsToPackage{unicode}{hyperref}
\PassOptionsToPackage{hyphens}{url}
%
\documentclass[
  11pt,
]{article}
\usepackage{lmodern}
\usepackage{amssymb,amsmath}
\usepackage{ifxetex,ifluatex}
\ifnum 0\ifxetex 1\fi\ifluatex 1\fi=0 % if pdftex
  \usepackage[T1]{fontenc}
  \usepackage[utf8]{inputenc}
  \usepackage{textcomp} % provide euro and other symbols
\else % if luatex or xetex
  \usepackage{unicode-math}
  \defaultfontfeatures{Scale=MatchLowercase}
  \defaultfontfeatures[\rmfamily]{Ligatures=TeX,Scale=1}
\fi
% Use upquote if available, for straight quotes in verbatim environments
\IfFileExists{upquote.sty}{\usepackage{upquote}}{}
\IfFileExists{microtype.sty}{% use microtype if available
  \usepackage[]{microtype}
  \UseMicrotypeSet[protrusion]{basicmath} % disable protrusion for tt fonts
}{}
\makeatletter
\@ifundefined{KOMAClassName}{% if non-KOMA class
  \IfFileExists{parskip.sty}{%
    \usepackage{parskip}
  }{% else
    \setlength{\parindent}{0pt}
    \setlength{\parskip}{6pt plus 2pt minus 1pt}}
}{% if KOMA class
  \KOMAoptions{parskip=half}}
\makeatother
\usepackage{xcolor}
\IfFileExists{xurl.sty}{\usepackage{xurl}}{} % add URL line breaks if available
\IfFileExists{bookmark.sty}{\usepackage{bookmark}}{\usepackage{hyperref}}
\hypersetup{
  hidelinks,
  pdfcreator={LaTeX via pandoc}}
\urlstyle{same} % disable monospaced font for URLs
\usepackage[margin=1in]{geometry}
\setlength{\emergencystretch}{3em} % prevent overfull lines
\providecommand{\tightlist}{%
  \setlength{\itemsep}{0pt}\setlength{\parskip}{0pt}}
\setcounter{secnumdepth}{-\maxdimen} % remove section numbering
\usepackage{mathrsfs}
\usepackage{esint}
\setmathfont{Asana Math}
\setlength{\emergencystretch}{3em}

\author{}
\date{23 September 2026}

\begin{document}

\hypertarget{ricci-energy-selection-scalar-coherence-and-einstein-defect}{%
\section{Ricci-energy selection, scalar coherence, and Einstein
defect}\label{ricci-energy-selection-scalar-coherence-and-einstein-defect}}

\hypertarget{the-question-this-note-resolves}{%
\subsection{1. The question this note
resolves}\label{the-question-this-note-resolves}}

The preceding Kähler calculation identified the limiting relative
Einstein defect of a constant-scalar-curvature blow-up sequence. Here we
remove the Kähler assumption and describe \textbf{every} sequence
approaching the sharp Ricci-energy bound. The remaining term is the
scalar variance. An actual conformal construction shows that this term
cannot be dropped, even on a fixed manifold admitting Kähler--Einstein
metrics and with zero primitive period cost. Conversely, choosing the
negative Yamabe representative removes this term without a spectral-gap
assumption.

The finite estimates below follow from classical Seiberg--Witten and
Chern--Weil inequalities. They do not require the project's integral
calibration lemma, its second-order estimates, or the Round21 spectral
criterion. This is a new deduction recorded in this research package,
not a claim of priority over the literature.

All manifolds are closed, connected, oriented, and four-dimensional.
Ricci norms are ordinary tensor norms; Weyl norms are endomorphism norms
on two-forms. Write \[
 T=2\chi+3\tau,\quad V=\operatorname{Vol}(g),\quad
 S_2=\int s^2,\quad R=\int|\operatorname{Ric}|^2,\quad
 P=\int|\operatorname{Ric}^{\circ}|^2,\quad
 L=\int(s-\bar s)^2.
 \tag{1}
\] Integrals use the current metric.

\hypertarget{precise-topological-input-and-the-finite-theorem}{%
\subsection{2. Precise topological input and the finite
theorem}\label{precise-topological-input-and-the-finite-theorem}}

Fix numbers \(A>0\), \(d\ge0\), with \(T=A-d\). For the metric in
question, assume an actual polarized monopole Spin\(^c\) structure whose
class \(a\) satisfies \(p=(a_g^+)^2\ge A\). Its equations must be
solvable throughout the conformal class, so that the classical scalar
and scalar--Weyl inequalities apply. A formal intersection vector is not
enough.

The class may depend on \(g\). No class-square upper bound, scalar band,
Sobolev constant, injectivity-radius bound, or diameter bound is
assumed.

In the standard application, \[
 M=X\#k\overline{\mathbb {CP}}^{\,2}
       \#m(S^1\times S^3),\qquad
 A=c_1^2(X)>0,\quad d=k+4m,
 \tag{2}
\] where \(X\) is a minimal complex surface of general type. The
polarized class selection and its \(b_+=1\) interpretation are as in
LeBrun's Proposition 3.2 and the
\href{lebrun_stability_closure.md}{application closure}. LeBrun's
Theorem 4.3 identifies the lower bound below as the actual all-metrics
infimum in this application.

Put \[
 C=8\pi^2A,\quad H=8\pi^2d,\quad
 J=R-(C+H).
 \tag{3}
\]

\textbf{Theorem 1 (finite energy-selection bound).} Under the stated
input, \(J\ge0\), and \[
 \boxed{\quad
 H-\sqrt{J(C+J)}\ \le P\le H+J .
 \quad} \tag{4}
\] The lower bound may of course be replaced by its positive part. With
\[
 \mathcal E(g)=\inf_{\lambda\ne0}
 \frac{\int|\operatorname{Ric}-\lambda g|^2}{\lambda^2 V},
 \tag{5}
\] one has the exact formula \[
 \mathcal E(g)=\frac{4P+L}{R}
 \tag{6}
\] and consequently \[
 \boxed{
 -\frac{4\sqrt{J(C+J)}}{C+H+J}
 \le
 \mathcal E(g)-\frac{4H+L}{C+H+J}
 \le\frac{4J}{C+H+J}.}
 \tag{7}
\]

Thus the Einstein defect is not determined by absolute Ricci energy
alone. Near its minimum, its two contributions are the topological
quantity \(4H\) and the scalar variance \(L\).

\hypertarget{proof}{%
\subsubsection{Proof}\label{proof}}

Define the classical LeBrun remainder \[
 D=\frac1{4\pi^2}\int
       \left(\frac{s^2}{24}+\frac12|W^+|^2\right)-\frac p2.
 \tag{8}
\] The required classical inputs are \[
 S_2\ge32\pi^2p,\qquad
 \int(-s+\sqrt6|W^+|)^2\ge72\pi^2p,\qquad D\ge0.
 \tag{9}
\] These are the scalar estimate, Theorem 2.3, and Proposition 4.1 of
\href{https://arxiv.org/pdf/math/0003068}{LeBrun, arXiv:math/0003068v1}.
No assertion about the equality set of a different curvature functional
is used.

Chern--Weil gives the exact identity \[
 R=8\pi^2(2p-T)+32\pi^2D,
 \qquad J=16\pi^2(p-A)+32\pi^2D.
 \tag{10}
\] In particular \[
 0\le p-A\le J/(16\pi^2),\qquad 0\le D\le J/(32\pi^2).
 \tag{11}
\] Let \[
 q=-s+\sqrt6|W^+|,\quad t=s+2\sqrt6|W^+|,\quad
 \delta=144\pi^2D,\quad
 \Delta=\int q^2-72\pi^2p,\quad U=\int t^2.
 \tag{12}
\] Here \(q\) need not be positive. The unsigned inequality in (9)
suffices. Pointwise expansion and integration yield \[
 \delta=\Delta+\tfrac12 U,\qquad \Delta\ge0,\qquad
 s=(t-2q)/3.
 \tag{13}
\] For the scalar surplus \(B_p=S_2-32\pi^2p\ge0\), this implies \[
 9B_p=4\delta-U-4\int qt.
 \tag{14}
\] Cauchy--Schwarz and \(U\le2\delta\) give \[
 B_p\le\frac49
       \left[\delta+\sqrt{2\delta(72\pi^2p+\delta)}\right].
 \tag{15}
\] Using (11), \(\delta\le9J/2\), and the convenient bound
\(72\pi^2p+\delta\le72\pi^2A+9J\), we obtain \[
 \boxed{0\le B_p\le2J+4\sqrt{J(C+J)}.}
 \tag{16}
\] There is room to sharpen these numerical constants; their optimality
is not asserted.

Since \(P=R-S_2/4\), (10) also gives \[
 P=H+J-8\pi^2(p-A)-B_p/4.
 \tag{17}
\] The upper bound in (4) follows immediately. For the lower bound, (11)
and (16) give \[
 P-H\ge J-J/2-\bigl(J/2+\sqrt{J(C+J)}\bigr)
          =-\sqrt{J(C+J)}.
 \tag{18}
\] Finally, expansion in \(1/\lambda\) yields \[
 \mathcal E=4-\frac{(\int s)^2}{VR}
            =\frac{4P+L}{R}.
 \tag{19}
\] Here \(R\ge C>0\). If \(\int s=0\), the infimum is \(4\), approached
as \(|\lambda|\to\infty\); otherwise the minimizer is
\(\lambda=R/\int s\). This covers both cases without falsely treating a
Ricci-flat metric as zero relative defect. Substitution proves (7).

\hypertarget{all-minimizing-sequences-not-only-kuxe4hler-sequences}{%
\subsection{3. All minimizing sequences, not only Kähler
sequences}\label{all-minimizing-sequences-not-only-kuxe4hler-sequences}}

\textbf{Corollary 2.} In the setting of Theorem 1, if \(J(g_j)\to0\),
then \[
 p_j\to A,\quad D_j\to0,\quad
 \int s_j^2\to32\pi^2A,\quad P_j\to8\pi^2d.
 \tag{20}
\] Moreover \[
 \boxed{\mathcal E(g_j)
   =\frac{4d}{A+d}
       +\frac{L_j}{8\pi^2(A+d)}+o(1).}
 \tag{21}
\] In particular \[
 \liminf\mathcal E(g_j)\ge\frac{4d}{A+d},
 \tag{22}
\] and \(\mathcal E(g_j)\to4d/(A+d)\) is equivalent to \(L_j\to0\). For
\(d=0\), vanishing relative Einstein defect is equivalent to vanishing
scalar variance along such a sequence.

Indeed (11) and (16) give (20). Since \(0\le L_j\le\int s_j^2\), the
variance is bounded; (7) therefore gives (21). This assertion imposes no
Kähler condition, no constant-scalar-curvature condition, and no
geometric compactness hypothesis.

For varying topologies the uniform formulation is (7). For example,
write \(r=d/A\), \(j=J/(8\pi^2A)\), and \(\ell=L/(8\pi^2A)\). Then the
difference between \(\mathcal E\) and \((4r+\ell)/(1+r+j)\) lies between
\(-4\sqrt{j(1+j)}/(1+r+j)\) and \(4j/(1+r+j)\). Small \emph{relative}
energy excess, rather than small unnormalized excess, is the appropriate
hypothesis when \(A\) varies.

\hypertarget{a-direct-scalar-coherence-criterion-on-the-original-metric}{%
\subsubsection{A direct scalar-coherence criterion on the original
metric}\label{a-direct-scalar-coherence-criterion-on-the-original-metric}}

For an ordinary monopole solution, put \(\rho=|\Phi|^2\) and let
\(\mathcal F\) be its real determinant curvature. The integrated spinor
identity and Hodge projection give \[
 B_p=\|s+\rho\|_2^2+8\|\nabla\Phi\|_2^2+
        8\|\mathcal F^+-(\mathcal F^+)_H\|_2^2.
 \tag{22a}
\] This retains the harmonic-projection loss in the classical scalar
inequality. It also follows from the \href{round17_integral.md}{Round17
identity} by subtracting the harmonic anti-self-dual contribution.

If \(s\ge-S\), the maximum principle yields \(0\le\rho\le S\). Since
\(|d\rho|^2\le4S|\nabla\Phi|^2\), Poincaré gives \[
 \sqrt L\le\|s+\rho\|_2+
       \sqrt{4S/\lambda_1}\,\|\nabla\Phi\|_2.
\] Joint Cauchy--Schwarz in the two displayed losses of (22a) proves \[
 \boxed{L\le\left(1+\frac{S}{2\lambda_1}\right)B_p.}
 \tag{22b}
\] Thus bounded \(S/\lambda_1\) is one sufficient condition for scalar
coherence along a minimizing sequence. It requires no weighted
equations, strict negative scalar band, or prior bound on Ricci error.
The stronger weighted variance estimate in the companion converse
theorem instead controls \(L\) linearly by \(D\), under its stated
negative-band assumptions. Neither argument supplies a spectral gap for
an arbitrary degenerating family.

\hypertarget{removing-scalar-variance-by-conformal-normalization}{%
\subsection{4. Removing scalar variance by conformal
normalization}\label{removing-scalar-variance-by-conformal-normalization}}

The variance in (21) is essential on the original metric. It can
nevertheless be removed by a mathematically legitimate change of
representative.

\textbf{Corollary 3 (Yamabe representatives of all minimizing
sequences).} In the standard application (2), or whenever Theorem 1's
input holds in the conformal class, let \(\widehat g\) be its negative
constant-scalar- curvature Yamabe representative. Then \[
 \widehat J\le J,\qquad
 R(g)-R(\widehat g)
       =\tfrac13\bigl(S_2(g)-S_2(\widehat g)\bigr),\qquad
 P(g)-P(\widehat g)
       =\tfrac1{12}\bigl(S_2(g)-S_2(\widehat g)\bigr).
 \tag{23}
\] Consequently every Ricci-energy-minimizing sequence can be replaced,
within its own conformal classes, by one satisfying \[
 \boxed{\mathcal E(\widehat g_j)\longrightarrow
                  \frac{4d}{A+d}.}
 \tag{24}
\] For \(d=0\) there is also the finite bound \[
 \boxed{\mathcal E(\widehat g)
          \le\frac{4J}{8\pi^2A+J}.}
 \tag{25}
\]

\hypertarget{proof-of-the-conformal-step}{%
\subsubsection{Proof of the conformal
step}\label{proof-of-the-conformal-step}}

A nonzero actual monopole projection rules out nonnegative scalar
curvature in this conformal class by the spinor Weitzenböck formula. The
Yamabe theorem therefore provides a representative with negative
constant scalar curvature. Four-dimensional Hodge star on two-forms,
\(p\), and \(\int|W^+|^2\) are conformally invariant.

For completeness, the scalar-energy minimization in a negative conformal
class has a direct verification. Rescale \(\widehat g\) to
\(s_{\widehat g}=-b<0\), and write \(g=u^2\widehat g\). With
\(\Delta=\operatorname{div}\nabla\), conformal covariance gives \[
 \int s_g^2\,d\mu_g
 =\int\left(-b-6u^{-1}\Delta_{\widehat g}u\right)^2d\mu_{\widehat g}.
 \tag{26}
\] Integration by parts yields \[
 \int u^{-1}\Delta u=\int u^{-2}|du|^2\ge0.
 \tag{27}
\] Expanding (26) shows \(S_2(g)\ge b^2V_{\widehat g}=S_2(\widehat g)\).
The Chern--Weil identity \[
 R=4\int|W^+|^2+\tfrac13 S_2-8\pi^2T
 \tag{28}
\] and \(P=R-S_2/4\) prove (23). The lower bound
\(R(\widehat g)\ge C+H\) remains valid for the same polarization and
actual class, so \(0\le\widehat J\le J\). Now (21) with \(\widehat L=0\)
proves (24). When \(d=0\), (4) gives \(\widehat P\le\widehat J\), whence
\[
 \mathcal E(\widehat g)=
 \frac{4\widehat P}{C+\widehat J}
 \le\frac{4\widehat J}{C+\widehat J}
 \le\frac{4J}{C+J}.
 \tag{29}
\]

This does not estimate the distance between \(g\) and \(\widehat g\).
The conformal scale-splitting example shows that such an estimate cannot
be inferred from \(J\to0\). Nor does an integral estimate for
\(\widehat g\) give a pointwise Ricci estimate.

\href{round14_global.md}{Round14} already carried out Yamabe
normalization for a particular shear construction. The assertion here
concerns arbitrary input conformal classes and all
Ricci-energy-minimizing sequences; it does not count the act of Yamabe
normalization as a new method.

\hypertarget{an-explicit-energy-exclusion-for-all-einstein-candidates}{%
\subsection{5. An explicit energy exclusion for all Einstein
candidates}\label{an-explicit-energy-exclusion-for-all-einstein-candidates}}

If a metric in this topological range is Einstein, then \(P=0\). Theorem
1 forces \[
 J(C+J)\ge H^2,\qquad
 \boxed{
 J\ge\frac{\sqrt{C^2+4H^2}-C}{2}
   =4\pi^2\bigl(\sqrt{A^2+4d^2}-A\bigr).}
 \tag{30}
\] For \(d>0\), any Einstein metric must therefore stay a positive
absolute-energy distance above the Ricci-energy infimum. This is a
consequence of the classical curvature inequalities, not a new
nonexistence theorem.

More generally, if \(0\le P/R\le\alpha\), then \[
 H\le\alpha(C+H+J)+\sqrt{J(C+J)}.
 \tag{31}
\] For the negative Yau26 normalization
\(\delta_-=\|\operatorname{Ric}^{\sharp}+I\|_{\infty,\mathrm{op}}<1\),
\[
 P\le4\delta_-^2V,\quad R\ge4(1-\delta_-)^2V,\quad
 \alpha=\frac{\delta_-^2}{(1-\delta_-)^2}
 \tag{32}
\] is admissible.

Thus, at any fixed positive blow-up/handle ratio \(d/A\), a pointwise
near-Einstein sequence cannot simultaneously be a relative Ricci-energy-
minimizing sequence. This is a usable restriction on a construction
strategy. It does not exclude inputs whose energy stays above that
minimum, and it does not decide whether the smooth manifold admits an
Einstein metric elsewhere.

\hypertarget{the-resulting-first-line-picture}{%
\subsection{6. The resulting first-line
picture}\label{the-resulting-first-line-picture}}

Three statements now have different, explicit roles.

\begin{enumerate}
\def\labelenumi{\arabic{enumi}.}
\tightlist
\item
  The existing integral gap near the forbidden \(2/3\) topology uses the
  project's integral calibration proof. It remains a theorem about
  exclusion of near-Einstein metrics in its actual-monopole range.
\item
  The present theorem describes every sequence minimizing the classical
  \emph{absolute} Ricci energy. It identifies the topological
  contribution to the relative defect and the independent
  scalar-coherence term.
\item
  The \href{line1_non_kahler_bridge.md}{fixed-metric converse} uses the
  earlier actual-form remainder and a scalar spectral gap. It starts
  from curvature-functional and period data, not from an assumed
  Ricci-small input. The \href{line1_conformal_neck.md}{scale-splitting
  example} explains why the scalar-coherence hypothesis has content.
\end{enumerate}

The prior actual-form supply uses \(t=1/2\), already an allowed negative
Kähler--Einstein root of Round21. Substituting that value into
\(\operatorname{dist}(t,\{0,1/2\})\) alone gives zero and proves no new
obstruction. The additional conclusions here come from curvature-energy
identities, actual topology, and control or removal of scalar variance.

The remaining Yau26 task is not concealed by these deductions: small
integrated Ricci error must still be connected to the required pointwise
input or to existence of an exact Einstein metric on the same smooth
manifold. Neither follows merely from minimizing the classical energy.
No new Yau26 branch is declared solved here.

\hypertarget{sources-and-proof-dependence}{%
\subsection{Sources and proof
dependence}\label{sources-and-proof-dependence}}

\begin{itemize}
\tightlist
\item
  LeBrun, \href{https://arxiv.org/pdf/math/0003068}{Ricci curvature and
  Seiberg--Witten theory}, cited preprint version: Theorem 2.3,
  Proposition 3.2, Propositions 4.1--4.2, and Theorem 4.3. These supply
  the inequalities, class selection, and sharp infimum.
\item
  \href{lebrun_stability_closure.md}{Polarized application closure}:
  project notation, metricwise class selection, and chamber details.
\item
  \href{line1_kahler_energy.md}{Previous Kähler calculation}: the
  Kähler-only minimizing-sequence profile, now extended above to all
  metrics with an explicit scalar-variance term.
\item
  \href{line1_scalar_band_remainder.md}{Actual-form remainder}: needed
  for the complementary fixed-metric bridge, not for Theorem 1.
\item
  The negative Yamabe representative uses the classical Yamabe theorem;
  all subsequent conformal monotonicity identities are proved in
  (26)--(29).
\end{itemize}

\end{document}
